%% kapa: exam=1 round=25 year=2014 period=1 form=A subject=civil_law subject_ko=민법 seq=7 number=7
%% source: Q-Net 감정평가사 기출문제 (cst00309) artlSeq=5101207 file=extracted/1008/5101207/2014년 제25회 감정평가사 1차시험 기출문제/1교시(A형).hwp
%% conversion: kordoc 4.24.51 -> tools/gen1.py (auto); review=proofread-sonnet; flags=none
\begin{question}{7}
민법 제35조에 의한 법인의 불법행위책임에 관한 설명으로 옳지 않은 것은? (다툼이 있으면 판례에 의함)
\begin{choices}
  \item 법령상 대표권 없는 이사의 행위에 대해서는 법인의 책임이 성립하지 않는다.
  \item 대표자로 등기되지 않았으나 법인을 실질적으로 운영하면서 법인을 사실상 대표하는 자의 직무수행에 대해서도 법인의 책임이 성립할 수 있다.
  \item 대표자 개인의 이익을 도모하기 위한 것이었더라도 외형상 객관적으로 법인의 대표자의 직무행위라고 인정할 수 있는 것이면, 법인의 책임이 성립할 수 있다.
  \item 대표기관의 고의적인 불법행위가 있으면 피해자에게 그 불법행위 내지 손해발생에 과실이 있더라도, 과실상계는 허용되지 않는다.
  \item 법인이 피해자에게 손해를 배상하면, 법인은 그 직무를 수행한 대표기관에 대해 선량한 관리자의 주의의무 위반을 이유로 구상권을 행사할 수 있다.
\end{choices}
\end{question}
\answer{4}
